\documentclass[12pt,a4paper]{article}

% ---------------------------------------------------------------
% METADADOS
% ---------------------------------------------------------------
\title{%
  Exist\^encia Global e Regularidade para Navier--Stokes Incompress\'ivel\\
  em $\mathbb{R}^3$: Barreiras Computacionais e Programa de Prova\\[4pt]
  {\large \textbf{Artigo F3} --- Alvo Clay --- \emph{Rascunho especulativo}}%
}
\author{%
  [Autor]\\
  {\small Federated Trivalent Manifold Research Programme}\\
  {\small \texttt{docs/pipeline/F3-Navier-Stokes/F3\_NS\_skeleton.tex}}
}
\date{2026 --- \textbf{NÃO SUBMETIDO} --- Lacunas abertas explícitas}

% ---------------------------------------------------------------
% PACOTES
% ---------------------------------------------------------------
\usepackage[T1]{fontenc}
\usepackage[utf8]{inputenc}
\usepackage[portuges]{babel}
\usepackage{amsmath,amssymb,amsthm}
\usepackage{hyperref}
\usepackage{geometry}
\usepackage{booktabs}
\usepackage{xcolor}
\usepackage{mdframed}
\geometry{margin=2.5cm}

% ---------------------------------------------------------------
% AMBIENTES
% ---------------------------------------------------------------
\newtheorem{theorem}{Teorema}
\newtheorem{lemma}[theorem]{Lema}
\newtheorem{proposition}[theorem]{Proposi\c{c}\~ao}
\newtheorem{corollary}[theorem]{Corol\'ario}
\theoremstyle{definition}
\newtheorem{definition}[theorem]{Defini\c{c}\~ao}
\newtheorem{remark}[theorem]{Observa\c{c}\~ao}

\newmdenv[backgroundcolor=green!10, linecolor=green!60!black,
          linewidth=1.5pt, roundcorner=4pt]{boxproved}
\newmdenv[backgroundcolor=red!8, linecolor=red!70!black,
          linewidth=1.5pt, roundcorner=4pt]{boxgap}
\newmdenv[backgroundcolor=blue!8, linecolor=blue!60!black,
          linewidth=1.5pt, roundcorner=4pt]{boxnumerical}
\newmdenv[backgroundcolor=orange!10, linecolor=orange!70!black,
          linewidth=2pt, roundcorner=4pt]{boxtodo}

% ---------------------------------------------------------------
\begin{document}
\maketitle

% ===== ALERTA CLAY ==============================================
\begin{mdframed}[backgroundcolor=red!12, linewidth=2.5pt, linecolor=red!70!black]
\textbf{\large ALERTA --- Este artigo NÃO está pronto para submissão ao Clay.}

\medskip
Regra Clay: publicado em Annals/JAMS/Inventiones/Acta, 2~anos sem furo,
comunidade aceita, Clay convoca comitê. Sem isso, não paga.

\medskip
\textbf{Buraco principal (TODO explícito):}
\[
  \texttt{TODO: provar } e_{\mathrm{linear}}(\omega) \ge 0
  \text{ sem \texttt{tanh}, para } N \to \infty.
\]
Enquanto $e_{\mathrm{lin}} = -0{,}046$ for instabilidade numérica
com regularização \texttt{tanh}, isto é \textbf{evidência numérica},
não blow-up analítico provado.

\medskip
Pré-requisito: os 24 artigos auditáveis (NS.A, NS.B, NS.C incluídos)
precisam ter DOI em mãos antes deste F3 ser submetido.
\end{mdframed}

\tableofcontents
\newpage

% ===== 1. INTRODUÇÃO ============================================
\section{Introdução e contexto}

O problema de Navier--Stokes (NS) do Clay pede:
dado dado inicial suave $u_0 \in C^\infty(\mathbb{R}^3)$ com
$\nabla \cdot u_0 = 0$ e $u_0$ de decaimento rápido,
determinar se existe solução suave global ou se ocorre
\emph{blow-up} em tempo finito.

Este artigo esqueleto descreve:
\begin{enumerate}
  \item As equações-alvo e obrigações formais herdadas de NS.A--NS.C;
  \item Os resultados já provados (barreira de regularidade);
  \item As evidências numéricas com parâmetros explícitos;
  \item Os \textbf{TODO} analíticos que precisam ser fechados;
  \item O que \emph{não} constitui prova Clay.
\end{enumerate}

% ===== 2. EQUAÇÕES-ALVO =========================================
\section{Equações-alvo e notação}

Para $u : \mathbb{R}^3 \times [0,\infty) \to \mathbb{R}^3$ incompressível:
\[
  \partial_t u + (u \cdot \nabla)u = -\nabla p + \nu \Delta u + f,
  \qquad \nabla \cdot u = 0.
\]

\textbf{Dado inicial:} $u_0 \in H^s(\mathbb{R}^3)$ com $s \ge \tfrac{5}{2}$,
$\|u_0\|_{H^s} = M < \infty$.

\textbf{Viscosidade:} $\nu > 0$ (caso físico relevante para Clay).

\textbf{Força:} $f \equiv 0$ (caso sem força, mais difícil).

\textbf{Normas críticas:} $\dot H^{1/2}$, $L^3$, $\mathrm{BMO}^{-1}$
(escala crítica por análise dimensional).

% ===== 3. RESULTADOS JÁ PROVADOS ================================
\section{Resultados já provados (não-Clay)}

\begin{boxproved}
\begin{theorem}[Existência local --- Leray 1934]
Para $u_0 \in H^1(\mathbb{R}^3)$, existe $T^* > 0$ e solução forte
$u \in C([0,T^*); H^1) \cap L^2(0,T^*; H^2)$ satisfazendo NS.
\end{theorem}
\textbf{Classificação: TEOREMA PUBLICADO.}
\end{boxproved}

\begin{boxproved}
\begin{theorem}[Solução fraca global --- Leray 1934]
Para $u_0 \in L^2(\mathbb{R}^3)$, existe solução fraca global
$u \in L^\infty(0,\infty; L^2) \cap L^2(0,\infty; \dot H^1)$
satisfazendo a desigualdade de energia de Leray.
\end{theorem}
\textbf{Classificação: TEOREMA PUBLICADO.}
Não resolve o problema Clay: falta unicidade e regularidade.
\end{boxproved}

\begin{boxproved}
\begin{theorem}[Critério de Serrin--Prodi--Ladyzhenskaya]
Se $u \in L^p(0,T; L^q)$ com $\frac{2}{p} + \frac{3}{q} \le 1$,
$q \ge 3$, então $u$ é regular em $[0,T]$.
\end{theorem}
\textbf{Classificação: TEOREMA PUBLICADO.}
Reduz o problema a: provar (ou refutar) que a solução de Leray
satisfaz o critério Serrin.
\end{boxproved}

\begin{boxproved}
\begin{theorem}[Estabilidade energética FTM --- NS2]
Na discretização FTM com $\nabla_h, \mathrm{div}_h$ compatíveis:
\[
  \|u_h(t)\|_2^2 + 2\nu \int_0^t \|\nabla_h u_h\|_2^2\,ds
  \le \|u_h(0)\|_2^2 + \text{forçamento}.
\]
\end{theorem}
\textbf{Classificação: TEOREMA CONDICIONAL}
(depende de compatibilidade do complexo discreto, verificada para
malhas regulares; caso anisotrópico: OPEN).
\end{boxproved}

% ===== 4. EVIDÊNCIAS NUMÉRICAS ==================================
\section{Evidências numéricas (não-Clay)}\label{sec:numericas}

\begin{boxnumerical}
\textbf{Resultado numérico NS-R1:}
Com $u_0 = A\,e^{-|x|^2}$, $\nu = 10^{-3}$, malha $N=256$,
regularização \texttt{tanh}:
\[
  e_{\mathrm{lin}} \approx -0{,}046 < 0.
\]
\textbf{Interpretação correta:} instabilidade linear numérica com
\texttt{tanh}. \textbf{NÃO é} blow-up analítico.\\
\textbf{Parâmetros completos:} Python 3.12, numpy 1.26, seed=2024,
dt=1e-4, T\_max=10.0, norma $H^1$ monitorada.
\end{boxnumerical}

\begin{boxnumerical}
\textbf{Resultado numérico NS-R2:}
Sem \texttt{tanh}: norma $H^1$ cresce em $t \in [0, 1{,}3]$,
atinge $\|u_h\|_{H^1} = 12{,}7$ (evidência de possível blow-up
numérico, sem garantia analítica).
\end{boxnumerical}

\begin{boxnumerical}
\textbf{Resultado numérico NS-R3:}
Vórtex de Taylor--Green com $\nu = 10^{-3}$: dissipação de energia
$\varepsilon(t)$ consistente com DNS, erro relativo $< 2\%$ até
$t = 5{,}0$ (validação do esquema, não prova de regularidade).
\end{boxnumerical}

% ===== 5. TODO ANALÍTICO ========================================
\section{TODO analítico --- Buracos para o Clay}\label{sec:todo}

\begin{boxtodo}
\textbf{TODO principal (F3-T1):}
\[
  \boxed{\text{Provar } e_{\mathrm{linear}}(\omega) \ge 0
         \text{ para } N \to \infty, \text{ sem \texttt{tanh}}.}
\]
Equivalentemente: mostrar que a linearização de NS em $u_0$ não
tem expoente de Lyapunov negativo no limite contínuo, ou construir
cenário de blow-up explícito satisfazendo as hipóteses Clay exatas.

\medskip
\textbf{Por que o dado atual não resolve:}
\begin{itemize}
  \item $e_{\mathrm{lin}} = -0{,}046$ usa \texttt{tanh} que
        regulariza artificialmente singularidades.
  \item Malha $N=256$ não captura escalas de Kolmogorov para
        $\nu \to 0$.
  \item Valor finito; limite $N \to \infty$ não controlado.
\end{itemize}

\textbf{Abordagem analítica necessária:}
Construir $u_0^*$ explícito tal que
$\limsup_{t \to T^*} \|\omega(\cdot,t)\|_{L^\infty} = \infty$
para algum $T^* < \infty$ (blow-up), OU provar estimativa
\emph{a priori}
\[
  \sup_{t \ge 0} \|u(\cdot,t)\|_{H^s} \le C(s,\nu,M) < \infty.
\]
sem usar compacidade que reintroduza solução fraca.
\end{boxtodo}

\begin{boxtodo}
\textbf{TODO (F3-T2): Estimativa crítica NS4}
\[
  \sup_{0 \le t < T^*} \|u(\cdot,t)\|_{\dot H^{1/2}} < \infty
\]
implica regularidade global (critério Escauriaza--Seregin--Sverak 2003,
caso $L^3$). Provar que solução de Leray satisfaz (ou viola) este
critério sem assumir resposta.
\end{boxtodo}

\begin{boxtodo}
\textbf{TODO (F3-T3): Compacidade e unicidade NS3}
Extrair limite $u_h \to u$ na topologia correta e mostrar que $u$
resolve NS no sentido forte (não apenas fraco).
Lacuna conhecida: termo convectivo $(u_h \cdot \nabla)u_h$ é
não-linear; compacidade forte de $u_h$ necessária mas difícil.
\end{boxtodo}

% ===== 6. O QUE NÃO CONSTITUI PROVA CLAY =======================
\section{O que \textbf{não} constitui prova Clay}

\begin{itemize}
  \item $e_{\mathrm{lin}} = -0{,}046$ com \texttt{tanh}:
        evidência numérica de instabilidade.
  \item Norma $H^1$ crescendo em simulação:
        evidência de possível singularidade numérica.
  \item Vórtex de Taylor--Green bem simulado:
        validação do esquema, não prova de regularidade global.
  \item Solução fraca de Leray existindo:
        provado em 1934; não resolve o Clay.
  \item Critério Serrin sendo satisfeito em região finita:
        verifica regularidade local, não global.
  \item Lean sem \texttt{sorry} para NS1--NS2:
        prova estabilidade energética discreta, não regularidade contínua.
\end{itemize}

% ===== 7. PROGRAMA MÍNIMO PARA SUBMISSÃO ======================
\section{Programa mínimo antes de submeter ao Clay}

\begin{enumerate}
  \item \textbf{DOIs dos 24 artigos} em mãos (NS.A, NS.B, NS.C
        publicados e citáveis).
  \item \textbf{F3-T1 fechado:} $e_{\mathrm{linear}}(\omega) \ge 0$
        analítico, sem \texttt{tanh}, sem $N$ finito.
  \item \textbf{F3-T2 fechado:} critério crítico $\dot H^{1/2}$
        satisfeito ou violado com construtividade.
  \item \textbf{F3-T3 fechado:} compacidade forte + unicidade.
  \item \textbf{Lean4 sem sorry} para todos os lemas analíticos.
  \item \textbf{Revisão bibliográfica} contra Constantin--Fefferman 1993,
        Caffarelli--Kohn--Nirenberg 1982, Tao 2016 (blow-up logarítmico).
  \item \textbf{Falsificador Popper:} qual dado inicial
        $u_0^*$ derruba a afirmação principal?
  \item \textbf{2 anos} em pre-print (arXiv) sem furo encontrado.
\end{enumerate}

% ===== 8. FALSIFICAÇÃO POPPER ==================================
\section{Falsificação Popper}\label{sec:falsificacao}

A afirmação principal deste programa é:

\begin{quote}
\emph{``O fluxo FTM-NS não desenvolve singularidade em tempo finito
para $u_0 \in H^s$ com $s \ge 5/2$, $\nu > 0$.''}
\end{quote}

\textbf{Falsificadores explícitos:}
\begin{enumerate}
  \item Construir $u_0^* \in C^\infty \cap L^2$ tal que
        $\|u(\cdot,T^*)\|_{\dot H^1} = +\infty$ para $T^* < \infty$.
  \item Mostrar que norma $H^s$ não pode ser uniformemente limitada
        para alguma família de dados iniciais grandes.
  \item Demonstrar que o esquema FTM viola a desigualdade de energia
        em alguma malha admissível.
\end{enumerate}

\textbf{Hipóteses ocultas a declarar:}
\begin{itemize}
  \item Simetria axial assumida em simulações não é hipótese Clay.
  \item Regularização \texttt{tanh} muda o sistema; remoção altera resultados.
  \item Dado inicial gaussiano não é o caso mais adverso.
\end{itemize}

% ===== 9. REFERÊNCIAS ===========================================
\section*{Referências}

\begin{enumerate}
  \item Leray, J. (1934). Sur le mouvement d'un liquide visqueux.
        \emph{Acta Math.}, 63, 193--248.
  \item Caffarelli, L., Kohn, R., Nirenberg, L. (1982).
        Partial regularity. \emph{Comm. Pure Appl. Math.}, 35, 771--831.
  \item Constantin, P., Fefferman, C. (1993). Direction of vorticity.
        \emph{Indiana Univ. Math. J.}
  \item Escauriaza, L., Seregin, G., Sverak, V. (2003). $L^3$ solutions.
        \emph{Russian Math. Surveys}, 58(2), 211--250.
  \item Fefferman, C. (2000). Navier-Stokes existence and smoothness.
        Clay Millennium Problem description. CMI.
  \item Tao, T. (2016). Finite time blowup for an averaged NS equation.
        \emph{J. Amer. Math. Soc.}, 29(3), 601--674.
  \item Clay Mathematics Institute (2000).
        \url{https://www.claymath.org/millennium-problems/}
\end{enumerate}

\end{document}
